\documentclass[UTF8,11pt]{ctexart}
\usepackage[a4paper,margin=24mm]{geometry}
\usepackage{amsmath,amssymb,amsthm,mathtools,mathrsfs}
\usepackage[all]{xy}
\usepackage{xurl,hyperref,enumitem}
\usepackage{tikz}
\usetikzlibrary{arrows.meta,cd,decorations.markings,calc,patterns}
\hypersetup{hidelinks,breaklinks=true}
\setlength{\emergencystretch}{3em}
\allowdisplaybreaks
\newtheorem*{theorem}{Theorem}
\newtheorem*{thm}{Theorem}
\newtheorem*{lemma}{Lemma}
\newtheorem*{proposition}{Proposition}
\newtheorem*{prop}{Proposition}
\newtheorem*{corollary}{Corollary}
\newtheorem*{cor}{Corollary}
\newtheorem*{claim}{Claim}
\theoremstyle{definition}
\newtheorem*{definition}{Definition}
\newtheorem*{defn}{Definition}
\newtheorem*{remark}{Remark}
\newtheorem*{example}{Example}
\newcommand{\R}{\mathbb R}
\newcommand{\C}{\mathbb C}
\newcommand{\Z}{\mathbb Z}
\newcommand{\Q}{\mathbb Q}
\newcommand{\N}{\mathbb N}
\newcommand{\F}{\mathbb F}
\newcommand{\D}{\mathbb D}
\newcommand{\bB}{\mathbb B}
\newcommand{\bC}{\mathbb C}
\newcommand{\bR}{\mathbb R}
\newcommand{\bZ}{\mathbb Z}
\newcommand{\bN}{\mathbb N}
\newcommand{\bQ}{\mathbb Q}
\newcommand{\bD}{\mathbb D}
\newcommand{\bF}{\mathbb F}
\newcommand{\CP}{\mathbb{CP}}
\newcommand{\RP}{\mathbb{RP}}
\newcommand{\sO}{\mathscr O}
\newcommand{\sI}{\mathscr I}
\newcommand{\sF}{\mathscr F}
\newcommand{\sG}{\mathscr G}
\newcommand{\sH}{\mathscr H}
\newcommand{\sR}{\mathscr R}
\newcommand{\sM}{\mathscr M}
\newcommand{\sV}{\mathscr V}
\newcommand{\sU}{\mathscr U}
\DeclareMathOperator{\Hom}{Hom}
\DeclareMathOperator{\Ext}{Ext}
\DeclareMathOperator{\Tor}{Tor}
\DeclareMathOperator{\Spec}{Spec}
\DeclareMathOperator{\Proj}{Proj}
\DeclareMathOperator{\supp}{supp}
\DeclareMathOperator{\im}{im}
\DeclareMathOperator{\id}{id}
\DeclareMathOperator{\Tr}{Tr}
\DeclareMathOperator{\tr}{tr}
\DeclareMathOperator{\rank}{rank}
\DeclareMathOperator{\ord}{ord}
\DeclareMathOperator{\dist}{dist}
\DeclareMathOperator{\End}{End}
\DeclareMathOperator{\Aut}{Aut}
\DeclareMathOperator{\Gal}{Gal}
\DeclareMathOperator{\vol}{vol}
\DeclareMathOperator{\Cl}{Cl}
\DeclareMathOperator{\coker}{coker}
\DeclareMathOperator{\codim}{codim}
\DeclareMathOperator{\divergence}{div}
\DeclareMathOperator{\Ric}{Ric}
\DeclareMathOperator{\grad}{grad}
\DeclareMathOperator{\Hess}{Hess}
\newcommand{\abs}[1]{\left\lvert#1\right\rvert}
\newcommand{\sabs}[1]{\lvert#1\rvert}
\newcommand{\babs}[1]{\bigl\lvert#1\bigr\rvert}
\newcommand{\Babs}[1]{\Bigl\lvert#1\Bigr\rvert}
\newcommand{\bbabs}[1]{\biggl\lvert#1\biggr\rvert}
\newcommand{\BBabs}[1]{\Biggl\lvert#1\Biggr\rvert}
\newcommand{\norm}[1]{\left\lVert#1\right\rVert}
\newcommand{\snorm}[1]{\lVert#1\rVert}
\newcommand{\bnorm}[1]{\bigl\lVert#1\bigr\rVert}
\newcommand{\Bnorm}[1]{\Bigl\lVert#1\Bigr\rVert}
\newcommand{\bbnorm}[1]{\biggl\lVert#1\biggr\rVert}
\newcommand{\BBnorm}[1]{\Biggl\lVert#1\Biggr\rVert}
\newcommand{\widebar}[1]{\overline{#1}}
\newcommand{\nicefrac}[2]{\frac{#1}{#2}}
\newcommand{\linnprod}[2]{\langle#1,#2\rangle}
\newcommand{\myindex}[1]{#1}
\newcommand{\glsadd}[1]{}
\newcommand{\avoidbreak}{}
\newcommand{\sourceref}[1]{\textnormal{[原文：\nolinkurl{#1}]}}
\newcommand{\thmref}[1]{\sourceref{#1}}
\newcommand{\lemmaref}[1]{\sourceref{#1}}
\newcommand{\propref}[1]{\sourceref{#1}}
\newcommand{\corref}[1]{\sourceref{#1}}
\newcommand{\defnref}[1]{\sourceref{#1}}
\newcommand{\exampleref}[1]{\sourceref{#1}}
\newcommand{\exerciseref}[1]{\sourceref{#1}}
\newcommand{\figureref}[1]{\sourceref{#1}}
\newcommand{\sectionref}[1]{\sourceref{#1}}
\newcommand{\appendixref}[1]{\sourceref{#1}}
\newcommand{\stacksref}[2]{\href{https://stacks.math.columbia.edu/tag/#1}{#2 [#1]}}


\newcommand{\E}{\mathbb E}
\newcommand{\Prob}{\mathbb P}
\newcommand{\Reff}{\mathcal R_{\mathrm{eff}}}
\newcommand{\eps}{\varepsilon}
\DeclareMathOperator{\diam}{diam}
\DeclareMathOperator{\Var}{Var}
\DeclareMathOperator{\Crit}{Crit}
\newcommand{\dd}{\,\mathrm d}


\begin{document}
\section*{083\quad 弱强制椭圆变分问题的Fredholm可解性与伴随障碍}
\subsection*{来源与整理范围}
Benoit Dionne, \emph{Partial Differential Equations}, \texttt{elliptic.tex}，Existence and Uniqueness of Solutions节，标签\texttt{ell\_exist\_th2}及紧邻的解算子/伴随构造，blob \nolinkurl{ad9fcaea78e33f90532a08f9bf5801832fa0076c}。
\par\url{https://github.com/BenoitDionne/PDE/blob/PDE/elliptic.tex}
\par\textbf{恢复方式（整理者说明）：}依据人类原文重新整理以补齐缺失题号；并非旧题证文件的逐字节恢复；2026-09-28。
\subsection*{作者命题与人类证明}

\paragraph{本节的环境与记号（据作者相邻段落整理）。}
$\Omega$为有界域，具有使$H^k(\Omega)\hookrightarrow L^2(\Omega)$紧的延拓/边界正则性；
$k\geq1$，$H^k_0(\Omega)\subseteq X\subseteq H^k(\Omega)$且$X$闭。
$\langle v,u\rangle_2=\int v\overline u$，$B:X\times X\to\mathbb C$为有界半双线性形式，线性于第一变量。
弱强制性为$\operatorname{Re}B(u,u)\geq C\|u\|_k^2-\lambda\|u\|_2^2$，$C,\lambda>0$。
\begin{theorem}[ell\_exist\_th2]
Let
\[
V=\{u\in X:B(v,u)=0\text{ for all }v\in X\},\qquad
W=\{w\in X:B(w,v)=0\text{ for all }v\in X\}.
\]
Then $\dim V=\dim W<\infty$. For $f\in L^2(\Omega)$, the variational
problem
\[
B(v,u)=\langle v,f\rangle_2\quad(v\in X),\qquad u\in X,
\]
has a solution if and only if $\langle w,f\rangle_2=0$ for all $w\in W$.
Solutions are unique up to addition of an element of $V$.
In particular, if $V=W=\{0\}$, each $f$ has a unique solution.
\end{theorem}
\paragraph{作者在主证明之前构造的解算子与伴随。}
If $B$ is coercive, Lax--Milgram provides a bounded linear operator
$F_B:L^2(\Omega)\to X$ satisfying
$B(v,F_Bf)=\langle v,f\rangle_2$. Here the functional is bounded because
$|\langle v,f\rangle_2|\leq\|f\|_2\|v\|_k$.
Define $B^*(u,v)=\overline{B(v,u)}$; it is likewise bounded and coercive,
so it has $F_{B^*}$. If $i:X\hookrightarrow L^2(\Omega)$ is the compact
inclusion, then $K_1=iF_B$ and $K_2=iF_{B^*}$ are compact, and
\begin{align*}
\langle K_1g,f\rangle_2
&=B^*(F_Bg,F_{B^*}f)
=\overline{B(F_{B^*}f,F_Bg)}\\
&=\overline{\langle F_{B^*}f,g\rangle_2}
=\langle g,K_2f\rangle_2.
\end{align*}
Thus $K_1^*=K_2$.
\begin{proof}[Proof of the theorem]
Choose $C>0,\lambda>0$ in the weak coercivity bound. Then
\[
\widetilde B(v,u)=B(v,u)+\lambda\langle v,u\rangle_2
\]
is coercive. By the preceding construction, now for $\widetilde B$,
there is a compact $K_1:L^2(\Omega)\to L^2(\Omega)$ with range in $X$
and
\[
\widetilde B(v,K_1f)=\langle v,f\rangle_2.
\]
For $u\in X$ we have the equivalences
\begin{align*}
 B(v,u)=\langle v,f\rangle_2\ (v\in X)
&\Longleftrightarrow\widetilde B(v,u)=\langle v,f+\lambda u\rangle_2\ (v\in X)\\
&\Longleftrightarrow u=K_1(f+\lambda u)\\
&\Longleftrightarrow \lambda^{-1}u-K_1u=\lambda^{-1}K_1f.
\end{align*}
Taking $f=0$ shows
\[
V=\ker(K_1-\lambda^{-1}I).
\]
Indeed an $L^2$ vector satisfying $u=\lambda K_1u$ automatically belongs
to $X$. The same construction for $B^*$ gives compact $K_2=K_1^*$ with
\[
W=\ker(K_2-\lambda^{-1}I).
\]
The compact-operator Fredholm theorem gives equal finite dimensions
of these two kernels, proving the first assertion.

The transformed equation is solvable exactly when
$K_1f\in\im(K_1-\lambda^{-1}I)$.
This range is closed by the same compact-operator theorem. The Hilbert
space range--kernel orthogonality identities give
\[
\im(K_1-\lambda^{-1}I)
=\ker(K_1^*-\lambda^{-1}I)^\perp=W^\perp.
\]
For $w\in W$,
\[
\langle K_1f,w\rangle_2
=\langle f,K_1^*w\rangle_2
=\lambda^{-1}\langle f,w\rangle_2.
\]
Therefore $K_1f\in W^\perp$ if and only if $f\in W^\perp$.
This is the solvability criterion. The difference of two solutions is
in $V$ by subtraction, and adding an element of $V$ preserves a solution.
The final assertion follows when the two kernels vanish.
\end{proof}

\subsection*{整理说明与先修范围（非作者正文）}
原章前文有有界域与紧嵌入约定，移到题面前显式保留，不能把此结论放到任意无界域。允许引用Lax--Milgram、Rellich紧嵌入以及紧算子的Fredholm备择；它们不是本题的待证结论。本题需构造带移位解算子，证明其伴随就是伴随形式的解算子，并把谱正交条件准确返回原右端项$f$，这些论证均保留。无关的例题和后续边界正则性未取入。独立计数目标是椭圆变分问题的核/余核与可解性，而不是把前题的紧嵌入或能量不等式另换名称。
\subsection*{可修改的简短标注（非作者正文）}
$c$：椭圆弱解的存在性及相容条件

$a$：弱强制性移位；Lax--Milgram解算子；紧嵌入；伴随；紧算子Fredholm备择；正交补

\texttt{cross\_theory\_status = yes}。作者把偏微分方程的半双线性形式转成恒等算子减紧算子的谱问题，通过显式伴随计算把有限维谱障碍识别为原方程伴随齐次解，最后恢复对$f$的相容条件。

全部标注为非穷尽、可修改初稿。
\subsection*{来源许可与编辑归属}
Benoit Dionne，\emph{Partial Differential Equations}，University of Ottawa，README署年2024，CC BY-NC-SA 4.0：\url{https://github.com/BenoitDionne/PDE/blob/PDE/README.md}；\url{https://creativecommons.org/licenses/by-nc-sa/4.0/}。本题编辑版本遵循同一许可。
ChatGPT 加入题号、中文来源说明、整理范围和初稿标注；编辑说明与人类证明分列。
\end{document}
